\section{Full GPT Results and Paired Analyses}\label{app:E}
\subsection{Allocation, configurations, and denominators}
The GPT record contains 688 allocated task stages: 288 v1 static tasks, 288 N2 static tasks, 48 CS1 static tasks, and 64 evidence-revision stages. The static total is therefore 624, not 688. Each static condition has R, A, and B. The eight recorded Codex configurations are 6 Astra, 6.1 Sol, 6.1 Terra, 6 Sol, 6 Luna, 5.6 Sol, 5.6 Terra, and 5.6 Luna, all with the recorded \texttt{xhigh} setting. These are eight product configurations from one provider family.

There are 684 standard-path raw files, of which 683 are nonempty. The frozen ledger classifies 651 stages as scored, seven as Schema failures, 25 as raw/parse failures, four as missing outputs, and one with a coarse output-limit label. The last label belongs to a CS1 \texttt{world:null} response and should not be read as a measured character overflow. The release retains the frozen status and this qualification. Appendix~\ref{app:J} reports failures separately from content errors.

Table~\ref{tab:E-main} uses every allocated determinate query opportunity. A failed output earns no correct readout but remains in the denominator. Table~\ref{tab:E-config} expands the result by configuration. The supplied B relations improve some constructions, but the improvement is neither universal nor sufficient to guarantee coherent geometry.

\begin{table}[htbp]\centering
\begin{tabular}{lrrr}\toprule Input & R & A & B\\\midrule
v1 & 599/600 & 541/600 & 550/600\\
N2 & 589/600 & 531/600 & 537/600\\
CS1 & 146/160 & 115/160 & 125/160\\\bottomrule
\end{tabular}
\caption{Correct determinate readouts under evaluator 2.2.1, including all allocated opportunities. v1 and N2 share buildings. CS1 B is 78.13\% under decimal half-up rounding.}\label{tab:E-main}
\end{table}

\begin{table}[htbp]\centering\small
\begin{tabular}{llrrrrr}
\toprule
Input & GPT config. & R & A & B & A valid & B valid \\
\midrule
v1 & 6 Astra & 75/75 & 71/75 & 75/75 & 12/12 & 12/12 \\
v1 & 6.1 Sol & 75/75 & 71/75 & 74/75 & 12/12 & 12/12 \\
v1 & 6.1 Terra & 75/75 & 68/75 & 74/75 & 11/12 & 12/12 \\
v1 & 6 Sol & 75/75 & 74/75 & 73/75 & 12/12 & 12/12 \\
v1 & 6 Luna & 74/75 & 63/75 & 63/75 & 11/12 & 11/12 \\
v1 & 5.6 Sol & 75/75 & 73/75 & 65/75 & 12/12 & 11/12 \\
v1 & 5.6 Terra & 75/75 & 66/75 & 61/75 & 11/12 & 10/12 \\
v1 & 5.6 Luna & 75/75 & 55/75 & 65/75 & 10/12 & 12/12 \\
N2 & 6 Astra & 74/75 & 63/75 & 69/75 & 11/12 & 11/12 \\
N2 & 6.1 Sol & 75/75 & 71/75 & 71/75 & 12/12 & 12/12 \\
N2 & 6.1 Terra & 74/75 & 75/75 & 75/75 & 12/12 & 12/12 \\
N2 & 6 Sol & 74/75 & 73/75 & 69/75 & 12/12 & 11/12 \\
N2 & 6 Luna & 72/75 & 55/75 & 69/75 & 10/12 & 12/12 \\
N2 & 5.6 Sol & 75/75 & 74/75 & 68/75 & 12/12 & 11/12 \\
N2 & 5.6 Terra & 72/75 & 53/75 & 58/75 & 10/12 & 10/12 \\
N2 & 5.6 Luna & 73/75 & 67/75 & 58/75 & 12/12 & 10/12 \\
CS1 & 6 Astra & 18/20 & 19/20 & 20/20 & 2/2 & 2/2 \\
CS1 & 6.1 Sol & 18/20 & 19/20 & 20/20 & 2/2 & 2/2 \\
CS1 & 6.1 Terra & 19/20 & 17/20 & 17/20 & 2/2 & 2/2 \\
CS1 & 6 Sol & 19/20 & 16/20 & 18/20 & 2/2 & 2/2 \\
CS1 & 6 Luna & 17/20 & 7/20 & 9/20 & 1/2 & 1/2 \\
CS1 & 5.6 Sol & 19/20 & 17/20 & 18/20 & 2/2 & 2/2 \\
CS1 & 5.6 Terra & 17/20 & 12/20 & 16/20 & 2/2 & 2/2 \\
CS1 & 5.6 Luna & 19/20 & 8/20 & 7/20 & 1/2 & 1/2 \\
\bottomrule
\end{tabular}

\caption{Full configuration results. v1/N2 each have 12 tasks per route, 75 determinate obligations per configuration; CS1 has two tasks and 20 obligations. A/B valid means contract-scorable, not geometrically correct.}\label{tab:E-config}
\end{table}

\subsection{A paired comparison on common-valid tasks}
The full denominator measures the delivered system, including format and delivery. To examine task-content differences, Table~\ref{tab:E-paired} pairs queries only when both compared outputs are contract-valid. It does not discard geometric errors. The resulting denominator can differ between R--A and A--B because validity depends on the compared routes.

\begin{table}[htbp]\centering\small
\begin{tabular}{llrrrrr}
\toprule
Input & Pair & $n$ & Both & First only & Second only & Neither \\
\midrule
N2 & R $\to$ A & 570 & 526 & 34 & 5 & 5 \\
N2 & A $\to$ B & 533 & 487 & 11 & 29 & 6 \\
CS1 & R $\to$ A & 140 & 113 & 15 & 2 & 10 \\
CS1 & A $\to$ B & 140 & 110 & 5 & 15 & 10 \\
\bottomrule
\end{tabular}

\caption{Common-valid paired outcomes. First-only means the first route is correct and the second is wrong or unevaluable. For R--A this is an answer-to-construction loss; for A--B, second-only is a relation-assisted rescue.}\label{tab:E-paired}
\end{table}

N2 has 34 R-correct/A-incorrect outcomes and five in the other direction among 570 common-valid obligations. The difference A minus R is $-5.09$ percentage points. A building-cluster bootstrap with 20,000 resamples and seed 20261105 gives a percentile interval of $[-9.29,-1.55]$ points. This interval describes the observed configuration ensemble on resampled buildings; it does not turn configurations or correlated questions into independent architectural samples.

For N2 A--B, B rescues 29 obligations and loses 11 among 533 common-valid pairs, a net difference of $+3.38$ points with interval $[0.19,6.42]$. Using all 600 allocated opportunities instead gives 50 rescues and 44 regressions, a net gain of six correct readouts. The two analyses answer different questions. The latter includes changes in output validity, while the former compares content on a shared valid subset. Neither should be substituted for the other without its denominator.

CS1 has 15 R-only versus two A-only outcomes, and 15 B-only versus five A-only outcomes, on 140 common-valid pairs for each contrast. Only two buildings contribute. We report these counts without a broad population confidence claim. The comparison is informative about these difficult drawings, but it is not a precise prevalence estimate for all complex staircases.

\subsection{Building and relation-type results}
Table~\ref{tab:E-building} shows the N2 and CS1 results by building, aggregated across the eight configurations. This aggregation keeps a hard building from disappearing in a high overall average. Table~\ref{tab:E-predicate} separates relation types. Both use full allocated opportunities; UNKNOWN is excluded from determinate accuracy and evaluated separately.

\begin{table}[htbp]\centering\small
\begin{tabular}{llrrr}
\toprule
Input & Building & R & A & B \\
\midrule
N2 & A01\_OCTAGON\_DC & 48/48 & 32/48 & 32/48 \\
N2 & A02\_DRAYTON & 45/48 & 45/48 & 48/48 \\
N2 & A06\_OCTAGON\_WI & 48/48 & 48/48 & 42/48 \\
N2 & B01\_nv0473 & 53/56 & 50/56 & 44/56 \\
N2 & B02\_mt0028 & 56/56 & 55/56 & 56/56 \\
N2 & B03\_az0061 & 47/48 & 44/48 & 48/48 \\
N2 & B04\_ca0801 & 39/40 & 36/40 & 39/40 \\
N2 & B05\_hi0032 & 48/48 & 36/48 & 37/48 \\
N2 & B06\_ga0030 & 48/48 & 48/48 & 48/48 \\
N2 & C01\_pr0070 & 48/48 & 42/48 & 36/48 \\
N2 & C02\_ut0027 & 46/48 & 41/48 & 47/48 \\
N2 & C03\_id0029 & 63/64 & 54/64 & 60/64 \\
CS1 & S01\_LOVELL\_HEALTH & 70/80 & 46/80 & 56/80 \\
CS1 & S02\_GAMBLE & 76/80 & 69/80 & 69/80 \\
\bottomrule
\end{tabular}

\caption{Correct readouts by building and route, pooling eight fixed configurations within each row. These are repeated measurements on a building, not independent buildings.}\label{tab:E-building}
\end{table}

\begin{table}[htbp]\centering\small
\begin{tabular}{llrrr}
\toprule
Input & Predicate & R & A & B \\
\midrule
N2 & above & 111/112 & 104/112 & 104/112 \\
N2 & connected\_by & 32/32 & 21/32 & 24/32 \\
N2 & disjoint & 56/56 & 54/56 & 53/56 \\
N2 & door\_host & 64/64 & 57/64 & 53/64 \\
N2 & identity\_assignment & 104/112 & 100/112 & 105/112 \\
N2 & same\_floor & 72/72 & 64/72 & 65/72 \\
N2 & stair\_connects & 15/16 & 12/16 & 11/16 \\
N2 & window\_host & 135/136 & 119/136 & 122/136 \\
CS1 & disjoint & 16/16 & 14/16 & 14/16 \\
CS1 & identity\_assignment & 59/64 & 49/64 & 56/64 \\
CS1 & same\_floor & 8/16 & 10/16 & 12/16 \\
CS1 & stair\_connects & 63/64 & 42/64 & 43/64 \\
\bottomrule
\end{tabular}

\caption{Determinate readouts by predicate. Differences combine relation difficulty, the available drawings, representation support, and delivery reliability.}\label{tab:E-predicate}
\end{table}

The granularity of a target can change the apparent error. Gamble Q01 is the ordinary \texttt{same\_floor} comparison, with R/A/B counts 8/8, 8/8, and 7/8. Gamble Q04 is a \texttt{stair\_connects} query, with 8/8, 4/8, and 2/8. Earlier narrative summaries incorrectly attached the ordinary-floor interpretation to Q04. The present tables use the frozen query IDs and predicates; no answer or score has been changed. Lovell Q04 is a separate local two-riser \texttt{same\_floor} question, with 0/8, 2/8, and 5/8. These examples show why a result must be traced to the actual query rather than to a remembered label.

\subsection{World coherence is a separate measurement}
Table~\ref{tab:E-world} adds geometry and coverage. All-relations-correct means that every determinate fixed query is correct for that world. Registered geometry checks the explicitly registered geometric roles and their dependencies. All-declared geometry also examines additional submitted entities. The joint result requires contract validity, required geometry coverage, correct determinate relations, and a PASS at the specified scope.

\begin{table}[htbp]\centering\small
\begin{tabular}{lrrrrrrr}
\toprule
Input & Alloc. & All rel. & Reg. & Joint reg. & Decl. & Joint decl. & Block \\
\midrule
v1/A & 96 & 77 & 46 & 43 & 9 & 9 & 5 \\
v1/B & 96 & 76 & 55 & 48 & 12 & 12 & 2 \\
N2/A & 96 & 69 & 50 & 42 & 16 & 14 & 1 \\
N2/B & 96 & 78 & 49 & 49 & 13 & 13 & 2 \\
CS1/A & 16 & 3 & 4 & 1 & 1 & 1 & 1 \\
CS1/B & 16 & 5 & 4 & 3 & 2 & 1 & 1 \\
\bottomrule
\end{tabular}

\caption{World-level results, out of allocated A/B worlds. Reg./Decl. are geometry PASS counts; joint also requires complete determinate correctness and coverage. Block means a definite represented material obstruction, a narrower class than all geometric violations.}\label{tab:E-world}
\end{table}

In N2, A has 69 all-relations-correct worlds and B has 78, but their all-declared joint counts are only 14 and 13 out of 96. A registered PASS is narrower: it does not certify every entity that the model chose to add. Conversely, an all-declared failure may concern a newly introduced fragment outside the registered geometric roles. This can be scientifically relevant without being a failure of a fixed local query.

N2 contains three definite material-blocking worlds. All three have every determinate fixed relation correct: 5.6 Luna's A and B Drayton Hall worlds, and its B Idaho City Schoolhouse world. In the latter, the internal staircase identity is itself queried by Q08, but its material geometry is outside the registered geometric-role list. It is therefore incorrect to call the internal stair wholly unqueried. The distinction is between an identity obligation and registered geometric scope.

\subsection{What the result establishes}
The observed configurations are generally strong at answering the selected relations, especially with semantic locators. Constructing a world imposes additional requirements that are not captured by a direct-answer score. Some losses are contract failures, some are interface or assembly errors, and a smaller set are explicit material conflicts. Relation assistance changes these outcomes but does not monotonically improve every building or every quality measure.

These results support measuring several output properties jointly. They do not establish a universal inability to reason spatially, a hidden internal world model, or a causal explanation based only on separate R/A generations. The casebook and common-valid analyses narrow the interpretation by identifying actual mechanisms and by retaining positive constructions as well as failures.
